\section{Part V: Discovering Customers in the Field}

\subsection{Interviews as Scientific Instruments}

Steve Blank's customer development methodology emphasizes leaving the building to test hypotheses with live market participants [R9]. An interview is not a sales pitch in disguise, nor a therapy session to validate the founder's ego. It is a \textbf{scientific instrument designed to falsify assumptions}.

Before scheduling a single conversation, formulate explicit prior hypotheses:
\begin{itemize}[leftmargin=1.2em,itemsep=0.2ex]
  \item Who experiences the acute struggle, and what triggered it?
  \item What specific makeshift workaround is currently used?
  \item What time and money are committed to existing remedies?
  \item Who holds financial authority, and who possesses veto power?
  \item What observable customer behavior would disprove our thesis?
\end{itemize}
Without prior hypotheses, founders instinctively cherry-pick pleasant compliments and mistake polite enthusiasm for validated demand.

\subsection{4 Fundamentally Different Research Conversations}

Do not use 1 generic script for every conversation. Deploy 4 specialized formats:
\begin{enumerate}[leftmargin=1.2em,itemsep=0.2ex]
  \item \textbf{Exploratory context interview:} Uncovers how the customer's daily operational world functions, where recurring bottlenecks arise, and how work is rewarded.
  \item \textbf{Decision reconstruction interview:} Forensically traces a recent, completed purchase or rejection, including rejected alternatives, internal debates, and budget approvals.
  \item \textbf{Contextual inquiry / workflow observation:} Directly observes the user performing the actual task in real time, noting unstated habits and manual workarounds [R11].
  \item \textbf{Offer evaluation interview:} Presents a precisely scoped, priced, bounded solution prototype to test trade-offs, pricing friction, and specific implementation objections.
\end{enumerate}

Never mix an exploratory interview with an aggressive sales pitch. Presenting a polished prototype prematurely contaminates the subject's memory and biases all subsequent answers.

\subsection{Who to Interview First: The Contrast Matrix}

When starting from zero, do not begin by interviewing personal friends, casual social media followers, or fellow founders. Construct a \textbf{Contrast Matrix}:
\begin{itemize}[leftmargin=1.2em,itemsep=0.2ex]
  \item \textbf{Recent buyers:} Individuals or managers who purchased an alternative in the past 6 months (reveals selection criteria and budget mechanisms).
  \item \textbf{Recent nonbuyers:} Individuals who thoroughly investigated solutions but decided to do nothing (reveals hidden friction and switching thresholds).
  \item \textbf{Active workaround builders:} Teams that built internal scripts or makeshift spreadsheets (reveals acute, unserved technical requirements).
  \item \textbf{Budget holders and approvers:} Finance managers, department deans, or engineering directors (reveals procurement hurdles and approval gates).
\end{itemize}

\subsection{Recruiting When You Begin with Zero Audience}

Cold outreach succeeds when framed as disciplined technical research rather than a disguised commercial pitch. Use legitimate professional networks, technical communities, and direct references.

A high-conversion research request must state:
\begin{itemize}[leftmargin=1.2em,itemsep=0.2ex]
  \item The precise technical situation you are studying (e.g., ``evaluating latency and hallucination trade-offs in production retrieval systems'').
  \item The exact qualifying criterion (e.g., ``engineers who shipped an AI feature in the last 6 months'').
  \item The explicit duration (``strictly 30 minutes, no sales pitch, no product demo'').
  \item Data protection and confidentiality commitments.
\end{itemize}
Never imply non-existent institutional partnerships, and never scrape private databases without consent.

\subsection{Sampling as a Research Design Problem}

There is no universal magic interview quota. 10 interviews in a tightly defined, homogeneous segment (e.g., backend engineers deploying production retrieval models) can reveal sharp, stable failure patterns. 50 interviews across a scattered, ill-defined group produce confusion.

Do not confuse qualitative mechanism discovery with quantitative market sizing. Jakob Nielsen's famous rule that 5 users discover 85\% of usability defects applies strictly to iterative usability testing of a concrete software interface, not to proving commercial demand or market size [R10]. Qualitative discovery uncovers \textbf{the causal mechanisms of choice}; quantitative surveys with rigorous sampling are required to estimate population prevalence and conversion rates [R12].

\subsection{Reconstruct Past Reality, Never Solicit Predictions}

The most fatal interview question is: ``Would you pay ₹40,000 for a live mentor-led AI course?'' It invites socially desirable flattery and speculative predictions. Human beings are notoriously inaccurate at forecasting their future self-discipline, budget allocation, and emotional stamina [R17].

Replace speculative questions with \textbf{forensic historical reconstruction}:
\begin{itemize}[leftmargin=1.2em,itemsep=0.2ex]
  \item Bad: ``Do you think continuous evaluation is important?''
  \item Good: ``Walk me through the exact steps you took the last time your model produced an embarrassing error in production. What tools did you open first?''
  \item Bad: ``Would your manager sponsor this training?''
  \item Good: ``What was the last external training course or tool your manager approved for your team? Who signed the invoice, and what documentation was required?''
\end{itemize}

\subsection{The 30-to-45-Minute Discovery Protocol}

A disciplined, non-leading discovery protocol follows 5 structured phases:
\begin{itemize}[leftmargin=1.2em,itemsep=0.2ex]
  \item \textbf{Phase 1: Context and Framing (3 to 5 mins):} Establish research goals, guarantee confidentiality, request recording permission, and explicitly instruct the subject: ``Please do not try to be polite; blunt criticism is 10 times more valuable to our research.''
  \item \textbf{Phase 2: The Triggering Incident (10 mins):} Anchor on a recent concrete episode. Establish date, project deadline, business stakes, and initial signs of failure.
  \item \textbf{Phase 3: The Workaround Journey (10 mins):} Trace every action taken: internal code, documentation consulted, free tools tested, hours invested, and results achieved.
  \item \textbf{Phase 4: Alternatives and Purchase History (10 mins):} Examine what external solutions were evaluated, who was consulted, what prices were reviewed, and what prevented adoption.
  \item \textbf{Phase 5: Synthesis and Referrals (5 mins):} Ask: ``What critical question did I fail to ask about this struggle? Who is 1 peer who experienced a completely different outcome in this situation?''
\end{itemize}

\subsection{Questions That Reveal Real Decisions}

Deploy high-resolution investigative prompts:
\begin{itemize}[leftmargin=1.2em,itemsep=0.2ex]
  \item ``What was happening in your company on the morning you decided you could no longer tolerate this problem?''
  \item ``How much time did you spend on this last week? Show me the calendar blocks or commit logs if accessible.''
  \item ``When you proposed buying a commercial solution, who pushed back, and what was their exact objection?''
  \item ``If this project had been canceled entirely, what would have been the financial or career consequence?''
\end{itemize}
Capture verbatim phrases. The specific nouns and verbs customers use to describe their agony form the core vocabulary of authentic product positioning.
